\color{unidarkgreen}{\subsection{Токовая защита нулевой последовательности}}
\color{black}

\begin{enumerate}

\item
Для батареи статических конденсаторов предусматривается два комплекта токовой защиты нулевой последовательности, предназначенных для резервирования действия основных быстродействующих защит при внутренних земляных видах коротких замыканий. Вывод защиты производится накладкой \textbf{«ТЗНП-1 ВВ»}.

\item
Измерительные цепи первого комплекта ТЗНП включены на токи фаз А, В, С со стороны высоковольтных вводов БСК. Защита реагирует расчетное значение утроенного тока нулевой последовательности, определяемого по формуле:

\begin{equation}
\label{eq:3I0}
3\dot{I}_{0} = \dot{I}_{\textsl{A}} + a^2 \cdot \dot{I}_{\textsl{B}} + a \cdot \dot{I}_{\textsl{C}}
\end{equation}

\item
Второй комплект ТЗНП может быть включен на токи фаз А, В, С со стороны нейтральных выводов БСК или к измерительному трансформатору току, установленному в цепи заземления нейтрали БСК. Данный комплект обладает высокой чувствительностью к однофазным коротким замыканиям вблизи нейтрали. 

\item
Оба комплекта выполнены идентично и содержат по две ступени, выполненных на базе модулей \textbf{«ТЗНП»}.  

\item
Каждая из ступеней защиты может выполнена с блокировкой по второй гармонике (отношению уровня второй гармоники к первой). Данная блокировка реализована на блоке  «ОВ БНТ ВВ». Выводится данная блокировка накладкой \textbf{«Режим БНТ ТЗНП-1 ВВ»}. 

\item
Защита имеет регулируемые значения по току и времени срабатывания, задаваемые параметрами \textbf{«3I0ср ТЗНП-1 ВВ»} и \textbf{«Тср ТЗНП-1 ВВ»} соответственно. 

\item 
Алгоритм работы ТЗНП НВ приведен на рисунке \ref{pict:TZNP_VV}. Алгоритм работы ТЗНП НВ выполнен аналогично. Уставки ТЗНП ВВ приведены в таблице \ref{app_set:tznpvv}. Уставки ТЗНП НВ приведены в таблице \ref{app_set:tznpnv}.

\medskip
\begin{figure}[H]
\centering
\includegraphics[width=1\textwidth,height=1\textheight,keepaspectratio]{\fbpath/09.01 БСК/271.0241 ТЗНП/img/TZNP_VV.pdf}
\captionof{figure}{Функциональная схема логики ТЗНП ВВ}\label{pict:TZNP_VV}
\end{figure}


\end{enumerate} 